% The class loads biblatex with backend=biber. This starter's compile
% pipeline runs the classic bibtex/(pdf|xe|lua)latex cycle (no biber pass), so
% the backend is switched to bibtex here, before the class loads biblatex —
% biblatex silently accepts a plain .bib file either way.
\PassOptionsToPackage{backend=bibtex}{biblatex}
\documentclass[9pt,a4paper,twocolumn,twoside]{rho-class/rho}
\usepackage[english]{babel}

%----------------------------------------------------------
% Title
%----------------------------------------------------------

\doctype{Research Article}
\title{A short, descriptive title for your paper}

%----------------------------------------------------------
% Authors, affiliations and dates
%----------------------------------------------------------

\author[{\textasteriskcentered},a,1]{Jane Smith}
\author[b,2]{Alex Chen}

\affil[a]{Department of Physics, Fictional State University}
\affil[b]{Institute for Applied Research, Fictional State University}

\dates{Manuscript compiled \today}

\corres{\textsuperscript{\textasteriskcentered}Corresponding author: \href{mailto:jane.smith@example.edu}{jane.smith@example.edu}}

\docinfo{This article was prepared with the rho-class LaTeX template and compiled with pdfLaTeX.}

\journalname{Working Papers}
\journal{Fictional State University}
\theday{\today}
\vol{1}
\no{1}

%----------------------------------------------------------
% Abstract and keywords
%----------------------------------------------------------

\begin{abstract}
    Replace this abstract with a concise, self-contained summary of the
    paper: the problem addressed, the method used, and the main result.
    Two to four sentences is typical for a research article.
\end{abstract}

\keywords{keyword one, keyword two, keyword three}

\begin{document}

\maketitle

\section{Introduction}

Start the introduction here. Motivate the problem, situate it against prior
work with a citation~\cite{example2024}, and state the contribution of this
paper in the last paragraph of the section.

\section{Background}

Use this section for the material a reader needs before the main
contribution: definitions, notation, or a summary of related methods.

\begin{note}
    The \texttt{note} environment is useful for calling out an assumption,
    caveat, or aside without breaking the flow of the main text.
\end{note}

\section{Methods}

Describe the approach taken. A typical equation is written as

\begin{equation}
    E = mc^2
    \label{eq:example}
\end{equation}

and referenced in text with \cref{eq:example}. The \texttt{stix2} font this
class loads covers the symbols needed for most equations without also
loading \texttt{amsmath}/\texttt{amssymb}.

\begin{info}
    Code listings can be added with the \texttt{listings} package, already
    loaded by this class:
\begin{lstlisting}[language=Python, caption={A minimal example.}]
def greet(name):
    return f"Hello, {name}!"
\end{lstlisting}
\end{info}

\section{Results}

Present the findings here. \Cref{fig:example} shows a placeholder figure —
replace \texttt{figures/example.pdf} with your own artwork.

\begin{figure}[htbp]
    \centering
    \fbox{\rule{0pt}{3cm}\rule{0.7\linewidth}{0pt}}
    \caption{Replace this placeholder with your own figure.}
    \label{fig:example}
\end{figure}

\section{Discussion}

Interpret the results, note limitations, and compare against the background
material introduced above.

\section{Conclusion}

Summarise the contribution and, if relevant, point to future work.

%----------------------------------------------------------
% Bibliography
%----------------------------------------------------------

\printbibliography

\end{document}
